% Titre sans numéro de chapitre, à la même place
\renewcommand{\chaptername}{\phantom{Chapitre}}\renewcommand{\thechapter}{}
\chapter{Les dictionnaires}
\renewcommand{\thesection}{\arabic{section}}
Dans ce chapitre nous présentons une nouvelle façon d'organiser et de traiter les données: \textbf{les dictionnaires}. Nous distinguerons la structure de données de son implémentation en illustrant les concepts au travers la réalisation de recettes de cuisine.\\[-0.5cm]

 La structure de donnée d'un dictionnaire est dite non linéaire.
\begin{tcolorbox}[enhanced,
    colback=green!25!black!10!white,colframe=green!75!black,title=Structures non linéaires,
    drop fuzzy shadow,watermark color=white]
Les données ne sont pas organisées dans un ordre, les éléments ne sont pas séquentiels, mais ils peuvent être reliés par une relation particulière.
  \end{tcolorbox}
Dans les cas des dictionnaires, les éléments sont reliés selon un couple appelé \textbf{clé, valeur}. Cette structure permet \textbf{d'accéder à un élément par sa clé} et non par un nombre (indice) comme les listes. Ci-dessous un exemple de dictionnaire où les clés sont les noms de fruits en Shimaoré et les valeurs leur nom français :
\begin{center}
    \noindent\begin{tabular}{l c l }
    "fenesi" &:& "jaque",\\
    "frampe" &:& "fruit a pain",\\
    "buyu" &:& "fruit du baobab"\\
    \end{tabular}
\end{center}
\vspace{0.5cm}
%
Quel est la valeur liée à la clé "frampe" ?
\vspace{0.5cm}

\noindent Comment faire une fonctionnalité équivalente avec deux listes ?
\vspace{2.3cm}


\subsection*{Comment les manipuler ? \textit{Interface}}

\noindent\begin{tabular}{l | l }
    \textbf{Ajouter} & ajoute un couple clé/valeur dans le dictionnaire.\\
    \textbf{Accéder à un élément} & accède à une valeur à partir de la clé.\\
    \textbf{Modifier un élément}  & modifie la valeur d'un élément\\
    \textbf{Supprimer un élément} & supprime un couple clé-valeur \\
    \textbf{Nombre d'éléments} & renvoie le nombre de couples contenus dans le dictionnaire. \\[0.5cm]
\end{tabular}

\subsection*{Comment les manipuler avec python ? \textit{Implémentation}}
En python, les dictionnaires sont contenus dans des accolades $\{ ... \}$.
Chaque clé est reliée à une valeur par deux points ($:$) et les couples clés/valeurs sont séparés par des virgules. Un exemple de dictionnaire où les clés sont les noms de fruits en Shimaoré et les valeurs leur équivalent en français:\\
\begin{center}
\noindent\begin{tabular}{l c l }
\{ & & \\
"fenesi" &:& "jaque",\\
"frampe" &:& "fruit a pain",\\
"buyu" &:& "fruit du baobab"\\
\} & & \\
\end{tabular}
\end{center}


\begin{lstlisting}[language=Python]
# Creer un dictionnaire vide 
d = {}

# Creer un dictionnaire non vide
d = {"fenesi" : "jaque", "frampe" : "fruit a pain", "buyu" : "fruit du baobab"}

# Taille du dictionnaire
print(len(d)) # affiche la taille du dictionnaire

# Acceder a un element par sa cle
print(d["buyu"]) # affiche "fruit du baobab"

# Ajouter un element
d["konokono"] = "corossol"
print(d) # affiche {"fenesi" : "jaque", "frampe" : "fruit a pain","buyu" : "fruit du baobab", "konokono" : "corossol" }

# Supprimer un element
del(d["frampe"])
print(d) # affiche {"fenesi" : "jaque", "buyu" : "fruit du baobab", "konokono" : "corossol" }
\end{lstlisting}

\vspace{0.3cm}
%
Voyons maintenant comment parcourir un dictionnaire. On imagine qu'on associe une quantité à chacun des fruits :
\begin{lstlisting}
MesIngredients = {"jaque" : 2, "fruit a pain" : 1, "fruit du baobab" : 2}
for cle, valeur in MesIngredients.items():
    print("j'ai ",valeur, " ", cle, " dans mon panier")
\end{lstlisting}
Que va afficher la console ? Expliquer la méthode items()
\vspace{1.7cm}

\noindent On peut également parcourir uniquement les clés du dictionnaires ou uniquement ses valeurs :
\begin{lstlisting}
MesIngredients = {"jaque" : 2, "fruit a pain" : 1, "fruit du baobab" : 2}
for cle in MesIngredients.keys():
    print(cle)
for valeur in MesIngredients.values():
    print(valeur)
\end{lstlisting}
Que va afficher la console ? Expliquer les méthodes keys() et values().
\vspace{1.7cm}

On ajoute à mon panier 2 fruits a pain et 1 jaque supplémentaire. Vous pouvez modifier votre liste de deux manières différentes :
\begin{lstlisting}
# On modifie la valeur en y accedant par la cle
MesIngredients["fruit a pain"] = 3 
# On incremente la valeur en y accedant par la cle
MesIngredients["jaque"] = MesIngredients["jaque"] + 1 
for cle, valeur in MesIngredients.items():
    print("j'ai ",valeur, " ", cle, " dans mon panier")
\end{lstlisting}
Que va afficher la console ?
\vspace{3cm}

On peut vérifier si une clé ou une valeur est dans un dictionnaire.
\begin{lstlisting}
# On verifie si une cle est dans MesIngredients
if "jaque" in MesIngredients :
    print(True)
else :
    print(False)
# On verifie si une valeur est dans MesIngredients
if 2 in MesIngredients.values() :
    print(True)
else :
    print(False)
if 3 in MesIngredients.values() :
    print(True)
else :
    print(False)
\end{lstlisting}
Que va afficher la console ?
\vspace{2cm}


\subsection*{Exercices}

\subsubsection*{Exercice 1}

Imaginons qu'on veut cuisiner un matsidza. On a déjà du riz et du curry et on va cueillir des noix de coco.
\begin{lstlisting}
MesIngredients = {
    "riz" : 1000,
    "curry" : True
} 
\end{lstlisting}

\begin{enumerate}[label=\alph*.]
    \item Ajoutez 3 noix de coco à vos ingrédients.
    \item Vous utilisez 200g de riz, 1 noix de coco et tout le curry. Modifiez vos ingrédients disponibles, un exemple est donné ci-dessous.
    \begin{lstlisting}
MesIngredients["riz"] = MesIngredients["riz"] - 200 \end{lstlisting}
    \item Affichez votre dictionnaire d'ingrédients dans la console.
\end{enumerate}

\subsubsection*{Exercice 2}
Nous créons deux dictionnaires qui contiennent les ingrédients nécessaires pour les deux recettes suivantes,  :
\begin{lstlisting}
pilao = {
    "riz" : 500,
    "poulets" : 500,
    "tomates" : 5,
    "boites de concentre de tomate" : 1,
    "oignons" : 5,
    "clous de girofle" : True,
    "sel" : True,
    "poivre" : True
}
matsidza = {"riz" : 200, "coco" : 1, "curry" : True}
\end{lstlisting}
\begin{enumerate}[label=\alph*.]
    \item Combien y a t-il de dictionnaires ?
    \item Comment s'appellent-ils ?
    \item Donner la taille de chacun des dictionnaires. Quelle fonction python utiliseriez-vous ?
    \item Quel est le type de la valeur reliée à la clé "sel" ?
\end{enumerate}


\subsubsection*{Exercice 3}
Nous créons trois autres dictionnaires :
\begin{lstlisting}
mataba = {
    "maniocs" : 5,
    "noix de coco" : 5,
    "oignons" : 2,
    "ail" : 2,
    "piment" : 1,
    "sel" : True
    }

mtsolola = {
    "bananes" : 5,
    "maniocs" : 1,
    "poissons" : 500,
    "oignons" : 2,
    "citrons" : 1, 
    "boite de tomate concentree" : 1,
    "cumin" : True,
    "poivre" : True,
    "sel" : True
    }

bata_bata = {
    "bananes" : 10,
    "maniocs" : 4,
    "sel" : True
    }
\end{lstlisting}
\begin{enumerate}
    \item Écrire un script qui parcourt les quatre recettes dictionnaires et qui affiche celles qui contiennent du riz
    \item Écrire un script qui parcourt les quatre recettes dictionnaires et qui affiche celles qui ont besoin de 5 bananes ou moins.
\end{enumerate}


\subsubsection*{Exercice 4 - à faire en binome}
On crée maintenant un dictionnaire pour répertorier les ingrédients disponible que nous avons.
\begin{lstlisting}
MesIngredients = {
    "poissons": 1000,
    "lait" : 1,
    "poulets": 1000,
    "oignons" : 10,
    "citrons" : 2,
    "boeuf" : 500,
    "boite de tomate concentree": 2,
    "bananes": 12,
    "ananas" : 3,
    "maniocs": 9,
    "songe" : 2,
    "tomates": 5,
    "riz": 10000,
    "noix de coco": 5,
    "ail": 2,
    "piment" : 2,
    "epices": ["safran","cumin","ail","poivre","sel","safran","cannelle","clou de girofle"]
}
\end{lstlisting}

Pour cette exercice vous pouvez le faire sur papier ou sur spyder au choix.

\begin{enumerate}
    \item Quelle est la particularité de la valeur liée à la clé "epices" ?
    \item Quelles recettes puis-je faire ?
    \item Ecrire un script qui vérifie votre réponse.
    \item Imaginons que vous avez executer la recette. Ecrire un script qui déduit de "MesIngredients" ceux que vous avez utilisés pour la/les recettes. \textit{Remarque : on a fini le sel mais pas les autres épices.}
    \item Comparez avec le reste de la classe. Avez vous fait la ou les mêmes recettes ?
\end{enumerate}





